\chapter{Ideen und Konzepte}
\label{cha:ideen_und_konzepte}

Dieses Kapitel begründet, wie die in Kapitel~\ref{cha:stand_der_technik} beschriebenen Bausteine Tor, Website-Fingerprinting, Circuit Padding und Shadow zu einem Versuchsaufbau zusammengeführt werden.
Die technische Umsetzung folgt in Kapitel~\ref{cha:realisierung}.

\section{Bewertung des Stands der Forschung}
\label{sec:konzepte_bewertung}

Padding und Website-Fingerprinting wurden in der Literatur untersucht, aber nicht in einer shadow Simulation systematisch gekoppelt.
Jansen und Wails messen \gls{df} in Shadow, variieren aber den Padding-Modus nicht~\cite{jansen2023data}.
Sie arbeiten zudem mit einem zellgenauen Cell-Trace aus einer gepatchten Tor-Version, während die vorliegende Arbeit am Monitor mitgeschnittene TCP-Pakete verwendet.
Die vorliegende Arbeit füllt diese Lücke, indem sie die drei in Tor angebotenen Padding-Konfigurationen unter sonst konstanten Bedingungen gegen \gls{df} misst.

\section{Konzeptioneller Versuchsaufbau}
\label{sec:konzepte_versuchsaufbau}

Der Versuchsaufbau folgt einem kontrollierten Laborexperiment, in dem der Padding-Modus gezielt variiert und alle übrigen Parameter konstant gehalten werden.
Die primäre Untersuchung erfolgt im Closed-World, ergänzt um ein einzelnes Open-World-Experiment als Realitätsanker.
Die formale Einordnung der Variablen und Messgrössen steht in Kapitel~\ref{cha:methoden}.

\section{Wahl des Simulators}
\label{sec:konzepte_simulator}

Die Umgebung muss vier Anforderungen erfüllen: realistische Verkehrsmuster, frei konfigurierbares Padding, exakte Reproduzierbarkeit und einen klar definierten Aufzeichnungspunkt.
Chutney und NetMirage sind primär als Tor-Test-Harness konzipiert und modellieren keinen realistischen Hintergrundverkehr~\cite{neverenough-sec2021}.
Live-Tor-Messungen sind nicht reproduzierbar~\cite{jansen2023data}.
Shadow mit Tornettools erfüllt alle Anforderungen: unmodifizierter Tor-Code und deterministischer Ablauf~\cite{netsim-atc2022}, und Jansen und Wails belegen mit 86\,\% Cross-Dataset-Accuracy die Eignung für \gls{wf}-Studien~\cite{jansen2023data}.

\section{Wahl des Klassifikators}
\label{sec:konzepte_klassifikator}

Aus den in Kapitel~\ref{cha:stand_der_technik} verglichenen Klassifikatoren wird \gls{df} gewählt.
\gls{df} arbeitet zeitunabhängig, weshalb simulationsbedingte Timing-Variationen die Ergebnisse nicht verfälschen.
Circuit Padding wirkt zudem direkt auf die von \gls{df} genutzten Richtungsfeatures, was den Padding-Effekt besonders sichtbar macht.
Die Validierung von \gls{df} in Shadow durch Jansen und Wails (82\,\% Accuracy auf Live-Tor) sichert die Vergleichbarkeit mit bestehender Literatur~\cite{jansen2023data}.

\section{Wahl des Szenarios}
\label{sec:konzepte_szenario}

Die primäre Untersuchung erfolgt im Closed-World, weil der Fokus auf dem isolierten Padding-Effekt liegt und nicht auf der Robustheit gegenüber unbekannten Seiten.
Die Closed-World-Accuracy ist zudem deutlich stabiler als die Open-World-Kennzahlen~\cite{jansen2023data}, was beobachtete Effekte klarer dem Padding zuordnen lässt.
Ergänzend wird ein Open-World-Experiment als Realitätsanker durchgeführt, ohne den Padding-Modus dort systematisch zu variieren.

\section{Wahl der Padding-Konfigurationen}
\label{sec:konzepte_padding_modi}

Der Padding-Modus nimmt drei Ausprägungen an: OFF, ON und Reduced.
Diese Wahl orientiert sich an den von Tor selbst angebotenen Optionen und erlaubt einen direkten Vergleich mit dem von Kadianakis et al.\ vorgeschlagenen Schema~\cite{kadianakis2021tor}.
Eine eigene Padding-Maschine wird nicht entwickelt, da das Ziel die Evaluation des produktiv eingesetzten Schemas ist und nicht der Entwurf einer neuen Verteidigung.

\section{Wahl der Tamaraw-Integration}
\label{sec:konzepte_tamaraw_integration}

Tamaraw ist nicht Bestandteil des regulären Tor-Daemons.
Im Gegensatz zum produktiv ausgerollten Circuit Padding plant das Tor-Projekt auch keine native Implementierung.
Tamaraw verzögert reale Zellen systematisch und erreicht in der ursprünglichen Auswertung rund 199\,\% Bandbreiten-Overhead bei zusätzlichem Latenz-Overhead~\cite{cai2014systematic}.
Beides widerspricht dem expliziten low-latency-Designziel von Tor, das eine Verzögerung realer Anwendungs-Zellen ausschliesst~\cite{kadianakis2021tor}.

Um Tamaraw dennoch in einem Tor-Kontext zu untersuchen, wird das Verfahren als externes Hilfsprogramm zwischen Monitor und Bridge realisiert, das den Zellfluss vor dem Versand abfängt, gemäss Tamaraw-Regelwerk formt und anschliessend wieder an Tor übergibt.
Die Einbindung erfolgt über Tors Pluggable-Transport-Schnittstelle\footnote{\url{https://spec.torproject.org/pt-spec/}}.
Ein \gls{pt} ist ein eigenständiges Hilfsprogramm, das der Tor-Daemon vor sich schaltet.
Jede Tor-Verbindung wird durch das Programm geleitet, das die Bytes vor dem Versand transformieren und beim Empfang zurücktransformieren kann.
Diese Schnittstelle ist in der Praxis vor allem zur Zensurumgehung (etwa obfs4) im Einsatz, ist aber generisch genug, um beliebige Strom-Transformationen einzuhängen, darunter auch eine WF-Verteidigung.

\section{Wahl des Hintergrundverkehrs}
\label{sec:konzepte_hintergrundverkehr}

Der Hintergrundverkehr stammt aus den von Tornettools bereitgestellten TGen-Clients, die das Anwendungsverhalten realer Tor-Nutzung über Markov-Modelle nachbilden~\cite{neverenough-sec2021}.
Eine eigene Verkehrsmodellierung wird nicht entwickelt, da die bestehenden Modelle bereits eine validierte Grundlage für realistischen Mischverkehr bilden.

\section{Wahl der Inhaltsquelle}
\label{sec:konzepte_inhaltsquelle}

Als Inhalt für die zu klassifizierenden Webseiten wird die Simple-English-Wikipedia verwendet.
Wikipedia liefert pro Artikel eine eigenständige HTML-Seite, was eine saubere Eins-zu-eins-Zuordnung zwischen Klasse und Inhalt ermöglicht. Da in Shadow kein Internetzugriff möglich ist wird ein lokaler Spiegel eingerichtet.


\section{Methodische Vereinfachung: Pcap statt Cell-Trace}
\label{sec:konzepte_pcap}

Jansen und Wails verwenden eine gepatchte Tor-Version, die pro Zelle einen beschrifteten Eintrag mit Circuit-Zuordnung in einen Cell-Trace schreibt~\cite{jansen2023data}.
Diese Variante ist nicht öffentlich verfügbar.

Stattdessen wird in dieser Arbeit am Monitor eine fortlaufende \gls{pcap} Datei mit allen TCP-Paketen aufgezeichnet und im Nachgang in beschriftete Richtungssequenzen pro Seitenbesuch zerlegt.

Während ein Cell-Trace die Pakete bereits zellgenau und mit Circuit-Information versehen liefert, ist das \gls{pcap} ein einziger unbeschrifteter Paketstrom, der erst nachträglich den einzelnen Besuchen zugeordnet werden muss.
Diese Zerlegung übernimmt ein eigens für die vorliegende Arbeit entwickeltes Python-Skript (vgl.\ Abschnitt~\ref{sec:realisierung_pcap_to_npz}), das die Pakete anhand der in der \texttt{schedule.json} protokollierten Besuchsfenster und Labels (Abschnitt~\ref{sec:konzepte_port_class}) den jeweiligen Seiten zuweist.

Damit arbeitet die vorliegende Pipeline auf TCP-Ebene und nimmt einen Detailverlust in Kauf.
Die zellgenaue Auflösung wird nicht erreicht, was gegenüber Cell-Traces einen Informationsverlust bedeutet.

Zugleich bildet dieses \gls{pcap}-Verfahren ein realistisches Angriffsszenario eines passiven Beobachters ab, dem ebenfalls keine internen Tor-Cell-Traces zur Verfügung stehen~\cite{wang2016realistically}.

Die Sequenzlänge wird nach Sirinam et al.\ auf 5000~Werte fixiert~\cite{sirinam2018deep}.
Jeder Wert entspricht in der vorliegenden Pipeline einem TCP-Paket innerhalb des Besuchsfensters, bei Sirinam et al.\ einer Tor-Zelle.
Aufzeichnungen mit mehr als 5000 Paketen werden am Ende abgeschnitten, kürzere Aufzeichnungen mit Nullen aufgefüllt.

\section{Merkmalswahl: nur Richtungen}
\label{sec:konzepte_features}

Als Eingangsmerkmal werden ausschliesslich die Paketrichtungen verwendet (ausgehend $+1$, eingehend $-1$), keine Zeitstempel.
Circuit Padding wirkt direkt auf die Richtungsabfolge, weshalb der Effekt im Richtungssignal am deutlichsten sichtbar wird.

\section{Zuordnung von Verkehr zu Webseiten}
\label{sec:konzepte_port_class}

Aus dem aufgezeichneten Paketstrom sollen beschriftete Trainingsbeispiele werden, jeweils bestehend aus einer Richtungssequenz und dem Label der abgerufenen Webseite.
Dazu muss jeder Seitenbesuch zweifach identifiziert werden: zeitlich, um die Besuche im fortlaufenden Paketstrom voneinander zu trennen, und inhaltlich, um die jeweils abgerufene Webseite zu benennen.

Die zeitliche Trennung folgt aus dem festen Visit-Intervall (vgl.\ Abschnitt~\ref{sec:konzepte_roundrobin}).
Jeder Monitor startet pro Intervall genau einen Seitenabruf, sodass jedem Zeitfenster genau ein Besuch entspricht.

Die inhaltliche Zuordnung ist aus dem Monitor-Pcap allein nicht rekonstruierbar.
Das Pcap zeichnet ausschliesslich den verschlüsselten Verkehr zum Tor-Guard auf, weder Ziel-IP noch Zielport des angesprochenen Webservers sind sichtbar.
Genau dieser Umstand entspricht dem realistischen Angreifermodell: dem passiven Beobachter zwischen Monitor und Guard fehlt eben jene Information.
Sie wird deshalb beim Aufbau der Simulation in einer Plandatei \texttt{schedule.json} (vgl.\ Abschnitt~\ref{sec:realisierung_schedule}) festgehalten, die pro Monitor und Besuch die Startzeit und die abgerufene Seite enthält.
Bei der Auswertung wird der \gls{pcap}-Strom an den Zeitgrenzen der Plandatei zerlegt und das hinterlegte Label übernommen.
Eine Markierung im Tor-Datenstrom selbst oder eine gepatchte Tor-Version, wie sie Jansen und Wails einsetzen~\cite{jansen2023data}, wird damit umgangen.


\section{Circuit-Isolation per NEWNYM}
\label{sec:konzepte_newnym}

Vor jedem Seitenbesuch wird über das Tor-Control-Interface ein neuer Circuit erzwungen.
Damit sind die einzelnen Besuche statistisch unabhängig und Verkehr aus früheren Besuchen kontaminiert die Zuordnung nicht.
Diese strikte Isolation ist eine Vereinfachung gegenüber dem realen Tor-Browser-Verhalten (First-Party-Bundling)~\cite{jansen2023data}, in der \gls{wf}-Literatur jedoch übliche Praxis.

\section{Round-Robin-Schedule}
\label{sec:konzepte_roundrobin}

Die Besuche werden klassenübergreifend nach einem Round-Robin-Verfahren verteilt: zuerst jede Seite einmal, dann jede ein zweites Mal und so fort.
Dadurch sind alle Klassen zeitlich gleichmässig über den Simulationszeitraum verteilt, was systematische Korrelationen zwischen Klasse und Netzwerkzustand (Last, Cache-Effekte) verhindert.
Zusätzlich bleibt die Klassen-Balance auch bei vorzeitigem Abbruch eines Laufs weitgehend erhalten.

\section{Skalierung der Topologie}
\label{sec:konzepte_skalierung}

Die Tor-Topologie wird auf 1\,\% des Live-Netzwerks skaliert, was rund 79~Relais ergibt~\cite{neverenough-sec2021}.
Eine Vollskalierung übersteigt den Arbeitsspeicher der bereitgestellten VM (128~GB) und macht die Simulationsdauer unpraktikabel für eine Bachelorarbeit.
Die Reduktion verringert die Pfaddiversität und macht die Verkehrsmuster für einen \gls{wf}-Klassifikator tendenziell leichter erlernbar~\cite{jansen2023data}, weshalb die gemessenen Accuracy-Werte als optimistische obere Schranke gegenüber dem Live-Netzwerk zu lesen sind.
Diese Limitation wird in Kapitel~\ref{cha:diskussion} wieder aufgegriffen.